\documentclass[10pt, letterpaper]{article}

% Packages:
\usepackage[
    ignoreheadfoot, % set margins without considering header and footer
    top=2 cm, % seperation between body and page edge from the top
    bottom=2 cm, % seperation between body and page edge from the bottom
    left=2 cm, % seperation between body and page edge from the left
    right=2 cm, % seperation between body and page edge from the right
    footskip=1.0 cm, % seperation between body and footer
    % showframe % for debugging 
]{geometry} % for adjusting page geometry
\usepackage{titlesec} % for customizing section titles
\usepackage{tabularx} % for making tables with fixed width columns
\usepackage{array} % tabularx requires this
\usepackage[dvipsnames]{xcolor} % for coloring text
\definecolor{primaryColor}{RGB}{0, 0, 0} % define primary color
\usepackage{enumitem} % for customizing lists
\usepackage{fontawesome5} % for using icons
\usepackage{amsmath} % for math
\usepackage[
    pdftitle={Hao Wang's CV},
    pdfauthor={Hao Wang},
    pdfcreator={LaTeX with RenderCV},
    colorlinks=true,
    urlcolor=primaryColor
]{hyperref} % for links, metadata and bookmarks
\usepackage[pscoord]{eso-pic} % for floating text on the page
\usepackage{calc} % for calculating lengths
\usepackage{bookmark} % for bookmarks
\usepackage{lastpage} % for getting the total number of pages
\usepackage{changepage} % for one column entries (adjustwidth environment)
\usepackage{paracol} % for two and three column entries
\usepackage{ifthen} % for conditional statements
\usepackage{needspace} % for avoiding page brake right after the section title
\usepackage{iftex} % check if engine is pdflatex, xetex or luatex

% Ensure that generate pdf is machine readable/ATS parsable:
\ifPDFTeX
    \input{glyphtounicode}
    \pdfgentounicode=1
    \usepackage[T1]{fontenc}
    \usepackage[utf8]{inputenc}
    \usepackage{lmodern}
\fi

\usepackage{charter}

% Some settings:
\raggedright
\AtBeginEnvironment{adjustwidth}{\partopsep0pt} % remove space before adjustwidth environment
\pagestyle{empty} % no header or footer
\setcounter{secnumdepth}{0} % no section numbering
\setlength{\parindent}{0pt} % no indentation
\setlength{\topskip}{0pt} % no top skip
\setlength{\columnsep}{0.15cm} % set column seperation
\pagenumbering{gobble} % no page numbering

\titleformat{\section}{\needspace{4\baselineskip}\bfseries\large}{}{0pt}{}[\vspace{1pt}\titlerule]

\titlespacing{\section}{
    % left space:
    -1pt
}{
    % top space:
    0.3 cm
}{
    % bottom space:
    0.2 cm
} % section title spacing

\renewcommand\labelitemi{$\vcenter{\hbox{\small$\bullet$}}$} % custom bullet points
\newenvironment{highlights}{
    \begin{itemize}[
        topsep=0.10 cm,
        parsep=0.10 cm,
        partopsep=0pt,
        itemsep=0pt,
        leftmargin=0 cm + 10pt
    ]
}{
    \end{itemize}
} % new environment for highlights


\newenvironment{highlightsforbulletentries}{
    \begin{itemize}[
        topsep=0.10 cm,
        parsep=0.10 cm,
        partopsep=0pt,
        itemsep=0pt,
        leftmargin=10pt
    ]
}{
    \end{itemize}
} % new environment for highlights for bullet entries

\newenvironment{onecolentry}{
    \begin{adjustwidth}{
        0 cm + 0.00001 cm
    }{
        0 cm + 0.00001 cm
    }
}{
    \end{adjustwidth}
} % new environment for one column entries

\newenvironment{twocolentry}[2][]{
    \onecolentry
    \def\secondColumn{#2}
    \setcolumnwidth{\fill, 4.5 cm}
    \begin{paracol}{2}
}{
    \switchcolumn \raggedleft \secondColumn
    \end{paracol}
    \endonecolentry
} % new environment for two column entries

\newenvironment{threecolentry}[3][]{
    \onecolentry
    \def\thirdColumn{#3}
    \setcolumnwidth{, \fill, 4.5 cm}
    \begin{paracol}{3}
    {\raggedright #2} \switchcolumn
}{
    \switchcolumn \raggedleft \thirdColumn
    \end{paracol}
    \endonecolentry
} % new environment for three column entries

\newenvironment{header}{
    \setlength{\topsep}{0pt}\par\kern\topsep\centering\linespread{1.5}
}{
    \par\kern\topsep
} % new environment for the header

\newcommand{\placelastupdatedtext}{% \placetextbox{<horizontal pos>}{<vertical pos>}{<stuff>}
  \AddToShipoutPictureFG*{% Add <stuff> to current page foreground
    \put(
        \LenToUnit{\paperwidth-2 cm-0 cm+0.05cm},
        \LenToUnit{\paperheight-1.0 cm}
    ){\vtop{{\null}\makebox[0pt][c]{
        \small\color{gray}\textit{Last updated in September 2024}\hspace{\widthof{Last updated in September 2024}}
    }}}%
  }%
}%

% save the original href command in a new command:
\let\hrefWithoutArrow\href

% new command for external links:


\begin{document}
    \newcommand{\AND}{\unskip
        \cleaders\copy\ANDbox\hskip\wd\ANDbox
        \ignorespaces
    }
    \newsavebox\ANDbox
    \sbox\ANDbox{$|$}

    \begin{header}
        \fontsize{25 pt}{25 pt}\selectfont Wang Hao

        \vspace{5 pt}

        \normalsize
        \mbox{Incoming Ph.D. Student, National University of Singapore (NUS)}%

        \vspace{3 pt}

        \normalsize
        \mbox{Zhejiang Hangzhou, China}%
        \kern 5.0 pt%
        \AND%
        \kern 5.0 pt%
        \mbox{\hrefWithoutArrow{3220104819@zju.edu.cn}{3220104819@zju.edu.cn}}%
        \kern 5.0 pt%
        \AND%
        \kern 5.0 pt%
        \mbox{\hrefWithoutArrow{tel:+86 152 7099 8779}{+86 152 7099 8779}}%
        \kern 5.0 pt%
    \end{header}

    \vspace{5 pt - 0.3 cm}

    \section{RESEARCH INTERESTS}
        \begin{onecolentry}
            \begin{highlightsforbulletentries}
                \item Structure-Preserving Operator Learning
                \item Scientific Machine Learning (SciML)
                \item Numerical Analysis for PDEs
                \item Computational Fluid Dynamics (CFD)
            \end{highlightsforbulletentries}
        \end{onecolentry}

    \section{EDUCATION}

        \begin{twocolentry}{
            Incoming
        }
            \textbf{National University of Singapore (NUS)}, Incoming Ph.D. Student
        \end{twocolentry}

        \vspace{0.10 cm}
        
        \begin{twocolentry}{
            Sep 2022 – Present
        }
            \textbf{Zhejiang University}, Mathematics and Applied Mathematics
        \end{twocolentry}

        \vspace{0.10 cm}
        \begin{onecolentry}  
            \begin{highlights}  
                \item GPA: 4.20/4.30 (93.70/100)
                \item \textbf{Relevant Coursework:} Real Analysis, Functional Analysis(H), PDE(H), Numerical Analysis.
            \end{highlights}  
        \end{onecolentry} 
    
    \section{PUBLICATIONS \& MANUSCRIPTS}
        \begin{onecolentry}
            \begin{highlights}
                \item \textbf{Hao Wang}, Kailiang Wu, et al. "	Learning Missing Physics from Legacy Simulators to Drive New Discovery from Model Error." \textit{Manuscript under submission}, 2025. \textbf{(First Author)}
                
                \item \textbf{Hao Wang}, Qi Tang. "Bridging Convection and Diffusion Regimes via a Structure-Preserving Neural Operator." \textit{Manuscript in preparation}, 2025.
            \end{highlights}
        \end{onecolentry}
    
    \section{RESEARCH EXPERIENCE}
        \begin{twocolentry}{
            Early 2025 – Present
        }
            \textbf{Remote Research Collaborator} \\
            \textit{Georgia Institute of Technology (Mentor: Prof. Qi Tang)}
        \end{twocolentry}

        \vspace{0.10 cm}
        \begin{onecolentry}
            \begin{highlights}
                \item Collaborating weekly on advanced topics in structure-preserving operator learning for problems in plasma and fusion modeling.
                \item Independently designed a novel solution embedding a semi-Lagrangian discretization into a Fourier Neural Operator to bridge advection and diffusion-dominated regimes.
                \item Developed and validated a single model that accurately captures both sharp shock structures and smooth convection-diffusion behavior.
            \end{highlights}
        \end{onecolentry}

        \vspace{0.2 cm} % Add some space between entries

        \begin{twocolentry}{
            Jul 2025 – Aug 2025
        }
            \textbf{Research Intern} \\
            \textit{Southern University of Science and Technology (Supervisor: Prof. Kailiang Wu)}
        \end{twocolentry}

        \vspace{0.10 cm}
        \begin{onecolentry}
            \begin{highlights}
                \item Single-handedly translated a conceptual framework (operator splitting with ML correction) into a mathematically rigorous and computationally robust model to bridge the model-reality gap.
                \item Conducted systematic analysis of the proposed operator-learning framework's stability and approximation properties.
                \item Implemented and benchmarked the framework on complex PDE systems, including Navier–Stokes, FitzHugh–Nagumo, and 1D compressible Euler equations.
            \end{highlights}
        \end{onecolentry}


    \section{RESEARCH PROJECTS}
        \begin{twocolentry}{
            Mar 2024 – May 2025
        }
            \textbf{Zhejiang Province SRTP Project:} Application and Optimization of Classical Computational Methods and Machine Learning Techniques in Flow Field Simulation \\ % Added a line break here if needed by your LaTeX setup
            \textit{Supervisor: Prof. Wang He-Yu} % Added supervisor line
        \end{twocolentry}

        \vspace{0.10 cm}
        \begin{onecolentry}
            \begin{highlights}
                \item Worked with a finite element method program provided by my instructor (Prof. Wang He-Yu) to solve the Navier-Stokes equations, gaining hands-on experience in computational fluid dynamics. % Optionally clarify here too
                \item Currently working on integrating classical computational methods with machine learning to enhance simulation accuracy and reduce computational costs. Exploring different optimization techniques and loss functions to improve the performance of the PINN model. 
            \end{highlights}
        \end{onecolentry}

    \section{HONORS \& AWARDS}
        \begin{onecolentry}
            \begin{highlights}
                \item \textbf{Chu Kochen Scholarship} (the highest honor bestowed upon ZJU students), Zhejiang University. 2025
                \item \textbf{National Scholarship} (Top 1\%), Awarded by the Ministry of Education, China. 2023, 2024, 2025 (Thrice)
                \item \textbf{Bronze Award} at the Shing-Tung Yau College Student Mathematics Competition (Applied and Computational Mathematics category), 2025
                \item \textbf{First-Class Scholarship}, Zhejiang University. 2023, 2024, 2025 (Thrice)
                \item \textbf{Finalist} at the Mathematical Contest in Modeling (MCM/ICM), 2024
            \end{highlights}
        \end{onecolentry}

    
    \section{SEMINARS \& READING GROUPS}
        \begin{onecolentry}
            \begin{highlights}
                \item \textbf{Fourier Analysis Reading Group}, Zhejiang University. Jul 2024  
                Studied \textit{Fourier Analysis: An Introduction} (E. M. Stein \& R. Shakarchi), covering fundamental concepts of Fourier series and transforms.

                \item \textbf{Finite Element Methods Reading Group}, Zhejiang University. 2024  
                Studied \textit{The Mathematical Theory of Finite Element Methods} (S. C. Brenner \& L. R. Scott), focusing on variational formulations, approximation theory, and error estimates.

                \item \textbf{Spectral Methods Group}, Zhejiang University. 2025  
                Studied \textit{Spectral Methods: Algorithms, Analysis and Applications} (Jie Shen, Tao Tang, Li-Lian Wang), emphasizing spectral discretizations for PDEs, algorithmic implementation, and convergence analysis.
            \end{highlights}
        \end{onecolentry}

    % NEW SECTION FOR THE WEBSITE INITIATIVE
    \section{LEADERSHIP \& INITIATIVES}
        \begin{twocolentry}{
             June 2024 – Present % e.g., Oct 2023 – Present
        }
            \textbf{Founder \& Maintainer: "Rhythm of Mathematics" Resource Sharing Website}
        \end{twocolentry}
        
        \vspace{0.10 cm}
        \begin{onecolentry}
            \begin{highlights}
                \item Developed and launched an online platform for Zhejiang University students, dedicated to sharing learning experiences and resources for various mathematics courses.
                \item Collaborated with the Zhejiang University Youth League Committee to promote and support the platform.
                \item Curated and organized a comprehensive collection of study materials, notes, and peer advice to foster a collaborative learning environment within the mathematics department.
            \end{highlights}
        \end{onecolentry}

    \section{TECHNICAL SKILLS}  
        \begin{onecolentry}  
            \begin{highlightsforbulletentries}  
                \item \textbf{Machine Learning:} Experienced in using PyTorch for setting up, training, and fine-tuning neural networks. Familiar with concepts such as convolutional neural networks (CNNs), Physics-Informed Neural Networks (PINNs).  
                \item \textbf{Operating Systems:} Proficient in Linux, with experience in using command-line tools for simulation setup, code execution, and data processing.  
            \end{highlightsforbulletentries}
        \end{onecolentry}  

\end{document}
